% !TeX program = lualatex
% !TeX encoding = UTF-8
% !TeX spellcheck = pt_BR

\documentclass[12pt,a4paper]{article}

\usepackage{fontspec}
\defaultfontfeatures{Ligatures=TeX}
\setmainfont{Latin Modern Roman}
\setsansfont{Latin Modern Sans}
\setmonofont{Latin Modern Mono}
\usepackage[brazilian]{babel}
\usepackage{geometry}
\geometry{top=2.2cm,bottom=2.2cm,left=1.5cm,right=1.5cm}
\usepackage{microtype}
\usepackage{ragged2e}
\usepackage{booktabs}
\usepackage{tabularx}
\usepackage{array}
\usepackage{enumitem}
\usepackage{xcolor}
\usepackage{listings}
\usepackage{xurl}
\usepackage{hyperref}
\usepackage{fancyhdr}
\usepackage{lastpage}

\definecolor{linkblue}{HTML}{1D4ED8}
\definecolor{codegray}{HTML}{F3F4F6}
\definecolor{rulegray}{HTML}{D1D5DB}
\definecolor{commentgreen}{HTML}{166534}
\definecolor{keywordblue}{HTML}{1D4ED8}
\definecolor{stringred}{HTML}{B91C1C}

\hypersetup{
  colorlinks=true,
  linkcolor=linkblue,
  urlcolor=linkblue,
  citecolor=linkblue,
  pdfauthor={Sergio Cordero Calvimontes},
  pdftitle={Relatório técnico do projeto arduino-uno-cooperative},
  pdfsubject={Concorrência cooperativa no Arduino Uno R3 / ATmega328P}
}

\pagestyle{fancy}
\fancyhf{}
\lhead{arduino-uno-cooperative}
\rhead{Relatório técnico}
\cfoot{Página \thepage\ de \pageref{LastPage}}
\renewcommand{\headrulewidth}{0.4pt}
\setlength{\headheight}{14.5pt}

\setlist[itemize]{leftmargin=1.5em,itemsep=0.25em,topsep=0.4em}
\setlist[enumerate]{leftmargin=1.7em,itemsep=0.25em,topsep=0.4em}
\setlength{\parindent}{1.2cm}
\setlength{\parskip}{0.25em}
\setlength{\emergencystretch}{2em}
\newcolumntype{Y}{>{\raggedright\arraybackslash}X}

\lstdefinestyle{cppstyle}{
  language=C++,
  basicstyle=\small\ttfamily,
  backgroundcolor=\color{codegray},
  frame=single,
  rulecolor=\color{rulegray},
  numbers=left,
  numberstyle=\scriptsize\color{gray},
  numbersep=8pt,
  keywordstyle=\color{keywordblue}\bfseries,
  commentstyle=\color{commentgreen},
  stringstyle=\color{stringred},
  showstringspaces=false,
  breaklines=true,
  columns=fullflexible,
  keepspaces=true,
  tabsize=2
}
\lstset{style=cppstyle}
\renewcommand{\lstlistingname}{Listagem}

\newcommand{\ProjectTitle}{Arduino Uno Cooperative}
\newcommand{\AuthorWebsite}{https://argus-kepheus.github.io/sergio-cordero-calvimontes/\#home}
\newcommand{\AuthorName}{\href{\AuthorWebsite}{Sergio Cordero Calvimontes}}
\newcommand{\RevisionDate}{19 de agosto de 2026}

\begin{document}

\begin{titlepage}
  \centering
  \vspace*{1.5cm}
  {\Large \textbf{Projeto pessoal}\par}
  \vspace{0.7cm}
  {\large Explorando concorrência cooperativa em um microcontrolador AVR de 8 bits\par}
  \vfill
  {\huge\bfseries \ProjectTitle\par}
  \vspace{0.8cm}
  {\Large Relatório técnico da arquitetura-base\par}
  \vfill
  \begin{tabular}{rl}
    Autor: & \AuthorName \\
    Plataforma de simulação: & Wokwi no navegador \\
    Linguagem: & C++ / Arduino \\
    Placa-alvo: & Arduino Uno R3 / ATmega328P \\
  \end{tabular}
  \vfill
  {\large \RevisionDate\par}
\end{titlepage}

\tableofcontents
\newpage

\section{Introdução}

O projeto \texttt{arduino-uno-cooperative} explora concorrência cooperativa em
um Arduino Uno R3 baseado no ATmega328P. A proposta deriva conceitualmente de
um projeto anterior em ESP32 com MicroPython e \texttt{asyncio}, mas não procura
traduzir literalmente aquela implementação. O objetivo é preservar as ideias de
múltiplas atividades lógicas, temporização não bloqueante, observabilidade e
degradação controlada sob restrições substancialmente mais severas de CPU, SRAM,
flash e quantidade de pinos.

A versão-base foi deliberadamente concebida sem RTOS, threads ou
\texttt{TaskScheduler}. Em seu lugar, utiliza um escalonador estático próprio,
com onze tarefas cooperativas. Seis dessas tarefas pertencem aos seis LEDs azuis,
mantidos independentes mesmo compartilhando o mesmo intervalo nominal.

\section{Objetivos técnicos}

Os objetivos principais são:

\begin{itemize}
  \item implementar onze tarefas cooperativas em um único fluxo físico;
  \item impedir espera bloqueante por meio da proibição de \texttt{delay()} na operação normal;
  \item observar tempo de callback, atraso, overruns, passagens e SRAM;
  \item operar simultaneamente LEDs, botões, serial, OLED e TFT;
  \item usar o próprio limite do ATmega328P como elemento do experimento;
  \item manter uma implementação-base legível para comparações futuras.
\end{itemize}

\section{Plataforma e distribuição de recursos}

A Tabela~\ref{tab:pinagem} apresenta a pinagem consolidada.

\begin{table}[htbp]
\centering
\caption{Pinagem oficial da arquitetura-base.}
\label{tab:pinagem}
\begin{tabularx}{\textwidth}{@{}lY@{}}
\toprule
Pino & Função \\
\midrule
D0/RX & recepção serial \\
D1/TX & transmissão serial \\
D2--D7 & seis LEDs azuis independentes \\
D8 & LED verde do botão principal \\
D9 & TFT D/C \\
D10 & TFT CS \\
D11 & TFT MOSI em SPI por software \\
D12 & LED laranja de atividade/ociosidade dos displays \\
D13 & TFT SCK em SPI por software e LED \texttt{L} integrado \\
A0 & botão principal \\
A1 & diminuir intervalo \\
A2 & aumentar intervalo \\
A3 & heartbeat amarelo do escalonador \\
A4/SDA & OLED SSD1306 \\
A5/SCL & OLED SSD1306 \\
\bottomrule
\end{tabularx}
\end{table}

A UART permanece preservada em D0/D1. O OLED utiliza o periférico I2C/TWI de
hardware. A TFT mantém fisicamente D11 como MOSI e D13 como SCK, porém utiliza o
caminho de SPI por software da biblioteca Adafruit. Essa decisão permite manter
D12 como GPIO de saída para o indicador laranja.

\section{Modelo de concorrência}

\subsection{Escalonador nativo}

Cada tarefa é descrita por um ponteiro de função, próximo instante de execução e
período. A estrutura é propositalmente pequena:

\begin{lstlisting}
struct Task {
    TaskCallback callback;
    TickMs nextRunMs;
    uint16_t periodMs;
};
\end{lstlisting}

Não existem pilhas privadas por tarefa. Uma tarefa executa até retornar; somente
então outra pode ser atendida. Por isso, callbacks longos são tratados como um
problema de projeto e devem ser decompostos em etapas.

\subsection{Relógio modular}

O escalonador utiliza os 16 bits inferiores de \texttt{millis()}. As comparações
são feitas por diferenças modulares assinadas. Como todos os períodos são muito
menores que 32768 ms, o sistema atravessa a transição \texttt{65535 -> 0} sem
precisar de tratamento especial de reset do relógio.

\subsection{Política de overrun}

Se um callback ou outra atividade atrasar o sistema a ponto de uma tarefa perder
uma ou mais liberações, essas liberações são descartadas. O escalonador não
executa múltiplos callbacks em rajada para recuperar o passado. A quantidade
descartada alimenta a métrica \texttt{OVR}.

\section{Tarefas da aplicação}

\begin{table}[htbp]
\centering
\caption{Tarefas da configuração-base.}
\begin{tabularx}{\textwidth}{@{}Y l Y@{}}
\toprule
Tarefa & Período & Responsabilidade \\
\midrule
\texttt{blinkLed1} a \texttt{blinkLed6} & 125--4000 ms & seis LEDs independentes \\
\texttt{scanButtons} & 5 ms & leitura e debounce \\
\texttt{sampleMetrics} & 250 ms & consolidar métricas \\
\texttt{serviceDisplays} & 20 ms & atualizar TFT/OLED em etapas \\
\texttt{printStatus} & 1000 ms & diagnóstico serial \\
\texttt{schedulerHeartbeat} & 100 ms & indicador amarelo em A3 \\
\bottomrule
\end{tabularx}
\end{table}

Os seis períodos dos LEDs são 125, 250, 500, 1000, 2000 e 4000 ms. O valor
inicial é 500 ms. A mudança de período procura preservar a posição relativa de
cada tarefa em seu ciclo atual.

\section{Entradas e debounce}

Os três botões são ativos em nível baixo e utilizam \texttt{INPUT\_PULLUP}. Uma
única tarefa de varredura atende os três botões. O debounce nominal é 30 ms,
com amostragem a cada 5 ms. O botão principal gera eventos de pressão e
liberação; os botões de intervalo reagem somente à pressão e não repetem a ação
quando permanecem pressionados.

\section{Subsistema de displays}

\subsection{TFT ILI9341}

A TFT é a interface principal e apresenta dois gráficos, campos numéricos e um
console circular de eventos. A aplicação não mantém framebuffer da TFT e evita
redesenhos completos. Cada gráfico avança apenas uma coluna por amostra, e
campos textuais possuem cache para que valores inalterados não sejam novamente
transmitidos.

\subsection{OLED SSD1306}

O OLED funciona como painel diagnóstico compacto. A biblioteca
\texttt{SSD1306Ascii} evita reservar o framebuffer de 1024 bytes que consumiria
metade da SRAM do ATmega328P. Sua atualização é limitada aproximadamente a 1 Hz.

\subsection{Pipeline cooperativo}

Cada amostra visual é congelada em \texttt{displaySnapshot}. A atualização é
então dividida em seis etapas: gráfico de \texttt{APP BUSY}, gráfico de SRAM,
três linhas de métricas TFT e OLED. Esse mecanismo reduz o maior bloco gráfico
executado em uma única passagem.

\section{Indicadores físicos}

O LED amarelo em A3 é alternado pela tarefa de heartbeat. O LED laranja em D12
fica em nível alto quando o subsistema de displays está logicamente ocioso e em
nível baixo durante operações instrumentadas. D13 também é conectado ao LED
\texttt{L} da própria placa, portanto atividade de clock da TFT pode aparecer
nesse LED integrado.

\section{Instrumentação}

A versão-base expõe:

\begin{description}
  \item[APP BUSY:] fração da janela de amostragem gasta dentro de callbacks instrumentados;
  \item[SRAM FREE:] estimativa instantânea da distância entre heap e pilha;
  \item[RAM LOW:] menor estimativa de SRAM livre amostrada desde a inicialização;
  \item[PASS/s:] passagens do escalonador por segundo;
  \item[MAX CALLBACK:] maior duração individual de callback;
  \item[MAX LATE:] maior atraso observado em relação à liberação programada;
  \item[OVR:] quantidade de liberações perdidas e descartadas.
\end{description}

\texttt{APP BUSY} não é apresentado como uso absoluto de CPU porque não inclui
perfeitamente todos os custos fora dos callbacks.

\section{Estratégia de memória}

A meta é manter pelo menos 512 bytes de SRAM livre estimada e não aceitar a
configuração integrada abaixo de 400 bytes sem nova revisão. A aplicação evita
\texttt{String}, \texttt{new}, \texttt{malloc}, containers dinâmicos e grandes
buffers. Estados pequenos são representados por máscaras e filas circulares de
tamanho fixo.

\section{Validação}

A pasta \texttt{tests/} contém sketches isolados destinados ao Wokwi web. Eles
não constituem testes automatizados do firmware integrado; servem para isolar
fiação, displays, botões e o escalonador. O arquivo
\texttt{docs/PT/validation-checklist.md} contém o roteiro de aceitação da versão
integrada.

Os cenários essenciais incluem operação a 125 ms com ambos os displays e serial
ativos, uso rápido dos três botões, geração de eventos suficiente para circular
o console, travessia do overflow modular e execução prolongada de pelo menos
15 minutos.

\section{Limitações e trabalhos futuros}

A versão-base não pretende ser a implementação de maior desempenho possível.
Permanecem abertos para trabalhos futuros: comparação com \texttt{TaskScheduler}
e AceRoutine, redistribuição de pinos para recuperar SPI de hardware, acesso
direto a registradores \texttt{PORT}, watchdog, prioridades, medição de jitter,
high-water mark real de pilha e persistência em EEPROM.

Manter esses itens fora da versão-base é intencional: o projeto precisa primeiro
fornecer um ponto de referência legível, instrumentado e reproduzível contra o
qual otimizações posteriores possam ser comparadas.

\section{Conclusão}

O \texttt{arduino-uno-cooperative} transforma as limitações do Arduino Uno em
parte do próprio experimento. Em vez de ocultar o custo das operações por meio de
um sistema operacional ou de abstrações de alto nível, a aplicação torna
explícitos o escalonamento, os atrasos, o consumo de SRAM e a degradação quando
prazos são perdidos. A arquitetura-base fica, assim, adequada tanto à simulação
no Wokwi quanto a futuras comparações de implementação no ATmega328P.

\end{document}
